% !TeX root = ../main.tex
\section{П.2 Характеристики шума на выходе демодулятора}

\vspace{1em}
\textbf{1. Расчет коэффициента передачи АД}

Согласно заданию, известны следующие параметры модели:
\begin{itemize}
    \item Эффективное напряжение шума на входе АД: $U_{\text{ш}} = 0.1 \text{ В}$
    \item Требуемое среднее значение выходного напряжения: $\bar{U} = 0.1 \text{ В}$
\end{itemize}

Для линейного амплитудного детектора (при воздействии только узкополосного нормального шума) среднее значение выходного напряжения определяется формулой математического ожидания распределения Релея:
\begin{equation*}
    \bar{U} = K_{\text{д}} U_{\text{ш}} \sqrt{\frac{\pi}{2}}
\end{equation*}

Выразим искомый коэффициент передачи АД ($K_{\text{д}}$):
\begin{equation*}
    K_{\text{д}} = \frac{\bar{U}}{U_{\text{ш}} \sqrt{\frac{\pi}{2}}}
\end{equation*}

Подставим заданные числовые значения:
\begin{equation*}
    K_{\text{д}} = \frac{0.1}{0.1 \cdot \sqrt{\frac{\pi}{2}}} = \frac{1}{\sqrt{\frac{\pi}{2}}} = \sqrt{\frac{2}{\pi}} \approx \sqrt{0.6366} \approx 0.798 \approx 0.8
\end{equation*}
Таким образом, для проведения моделирования необходимо задать значение коэффициента передачи детектора $K_{\text{д}} = 0.8$.

\vspace{1em}
\textbf{2. Теоретические значения для сопоставления}

Для выполнения второй части задания (сопоставление измеренных значений с теоретическими), зафиксируем расчетные показатели:
\begin{itemize}
    \item \textbf{Теоретическое среднее напряжение:} $\bar{U}_{\text{теор}} = 0.1 \text{ В}$ (задано по условию).
    \item \textbf{Теоретическое СКО (эффективное напряжение) на выходе АД:}
\end{itemize}
Рассчитаем СКО флуктуаций $\sigma_U$ на выходе АД по стандартной формуле для дисперсии огибающей:
\begin{equation*}
    \sigma_U = K_{\text{д}} U_{\text{ш}} \sqrt{\frac{4 - \pi}{2}}
\end{equation*}

Используя точное значение $K_{\text{д}} = \sqrt{2/\pi}$:
\begin{equation*}
    \sigma_{U\text{ теор}} = \sqrt{\frac{2}{\pi}} \cdot 0.1 \cdot \sqrt{\frac{4 - \pi}{2}} = 0.1 \cdot \sqrt{\frac{4 - \pi}{\pi}} = 0.1 \cdot \sqrt{\frac{4}{\pi} - 1} \approx 0.1 \cdot 0.5227 \approx 0.0523 \text{ В} = 52.3 \text{ мВ}
\end{equation*}

\textbf{Вывод:} При проведении моделирования с заданным коэффициентом $K_{\text{д}} = 0.8$, измерительные приборы должны показать среднее напряжение около $0.1 \text{ В}$ и эффективное напряжение (СКО) около $52.3 \text{ мВ}$. Эти теоретические значения следует занести в отчет для сопоставления с результатами измерений.